% Source: sources/returns/S6_04_MAPS_MANIFOLDS/sections/manifolds.tex:486–613
\paragraph{Theorem D.1.9 (Proved: the finite-coefficient cubic).}
Let $d>1$ and let $C_\delta$ be the complex-oriented fixed curve of
the $\delta$ action downstairs. There is a class
$v\in H^2(X_p;\mathbb Z/d)$ with $v(C_\delta)=1$.
For $u=\pi^*v$ and the reductions of the classes $x_j$ in D.1.6,
the full degree-two cubic, evaluated on the complex fundamental
class of $\widetilde X_p$, is
\[
 \begin{aligned}
 \langle x_i x_j x_k,[\widetilde X_p]\rangle&=0,
 &\langle u x_i x_j,[\widetilde X_p]\rangle&=C_{ij},\\
 \langle u^2x_i,[\widetilde X_p]\rangle&=-12,
 &\langle u^3,[\widetilde X_p]\rangle&=0
 \quad\text{in }\mathbb Z/d.
 \end{aligned}
\]
where $C_{ii}=-2$, $C_{i,i\pm1}=1$ and all other entries are zero.
The $d-1$ classes $x_0,\ldots,x_{d-2}$ together with $u$ form a
$(\mathbb Z/d)$-basis of $H^2(\widetilde X_p;\mathbb Z/d)$.
In particular $u$ is deck invariant. No nontrivial extension of the
deck representation is needed to compute this cubic.
\begin{proof}
The fixed curve generates $H_2(X_p)=\mathbb Z/d$ by D.1.6's
downstairs calculation. The universal coefficient epimorphism
\[
 H^2(X_p;\mathbb Z/d)\longrightarrow
 \operatorname{Hom}(H_2(X_p),\mathbb Z/d)
\]
therefore supplies $v$. Each lift of $C_\delta$ maps to it with
degree one, so $\langle u\kappa,[\widetilde X_p]\rangle=1$.
The diagonal-direction external curve has zero homology class,
and the ruling is its sum with the fixed curve. Thus the assertion
about this pairing uses the complex ruling orientation, not a
choice of a torsion generator up to sign. Universal coefficients
upstairs identify degree-two cohomology with
$\operatorname{Hom}(\mathbb Z^{d-1}\oplus\mathbb Z/d,\mathbb Z/d)$.
Evaluation on its torsion summand shows that $u$ completes the
reductions of the free integral basis to a basis. We choose free
lifts in $H_2$ annihilated by $u$, subtracting multiples of its
torsion generator if necessary; this fixes the splitting used in
the classifying coordinates below. The first two
cubic formulas follow from D.1.6.

We compute the remaining component restriction with its orientation.
In the toric model with lattice coordinates $(r,s)$ on its height-one
vertices, take the integral support function
\[
 \psi(r,s)=-\frac{r(r-1)}2+3s(s-1),
\]
interpolated linearly over each unimodular cone. It defines a toric
line bundle. Its changes under the two period translations are
\[
 \psi(r+1,s)-\psi(r,s)=-r,
 \qquad
 \psi(r,s-1)-\psi(r,s)=-6s+6.
\]
These affine functions give commuting character multipliers
$x_1^{-1}$ and $x_2^{-6}t^6$, so the bundle descends over the
period lattice generated by $e_1$ and $-e_2$. The holomorphic-unit
corrections in the actual cusp action can be deformed to one in a
sufficiently small tube, as in the checked cusp isotopy of AC26,
Section~7.2. This transports the integral topological class and
does not change any of the following pairings.

For a wall with endpoints $v,v'$ and the two adjacent third vertices
$w,w'$, the relation of height-one rays is $w+w'=v+v'$.
The two toric trivializations on the curve differ by a power of its
coordinate of exponent
$\psi(v)+\psi(v')-\psi(w)-\psi(w')$. This is its degree for the
complex orientation. Using respectively
\[
 \begin{array}{c|c|c}
 (v,v')&(w,w')&\deg\\\hline
 ((0,0),(1,0))&((0,1),(1,-1))&-6\\
 ((1,0),(0,1))&((0,0),(1,1))&0\\
 ((0,0),(0,1))&((1,0),(-1,1))&1
 \end{array}
\]
gives the internal, diagonal and $\delta$-fixed curve degrees
$(-6,0,1)$. The regular-fibre first Chern class of this bundle is
$-(u\wedge w+6\gamma\wedge\delta)$: the winding of the first
multiplier on the ordered $(\widehat u,\widehat w)$ rectangle is
$-1$, and that of the second on
$(\widehat\gamma,\widehat\delta)$ is $-6$.
There are no other windings. This also fixes the relative signs
that AC26 (7.9), when used with separately oriented curve generators,
leaves implicit.

The complement class dual to $\Phi=(4,2,3,2)$ restricts to twice
$q=u\wedge w+6\gamma\wedge\delta$. The specialization from the
regular fibre to the downstairs cusp is onto in degree two
(AC26, Proposition~7.12, whose integral annulus proof is used here).
Explicitly, for any of the three complex curves $C$, choose
$y\in H_2(F;\mathbb Z)$ specializing to $[C]$. The descended
toric line bundle evaluates on $y$ as $-q(y)=\deg_L(C)$,
while its image in the exterior cyclic generator $g$ has
coefficient $2q(y)$. Thus, integrally modulo $d$,
\[
 [C]=-2\deg_L(C)\,g\quad\text{in }H_2(X_p;\mathbb Z).
\]
Since $2$ is a unit modulo $d$, the images of the three
complex curves have precisely the ratios $-6:0:1$ after
normalizing by the complex $\delta$-fixed curve. On the normalized $dP_6$ the
internal curve is $E_1$, the zero external curve is $E_2$, and the
fixed external curve is $E_3$. The opposite-pair relation is
$H=E_1+E_2+E_3$ after mapping to the cusp. Thus the restriction of
$u$ to every normalized component is
\[
 \eta=-5H+6E_1-E_3=-5(H-E_1)+(E_1-E_3)
 \quad\text{in }H^2(dP_6;\mathbb Z/d).
\]
Indeed its evaluations on $E_1,E_2,E_3$ are $-6,0,1$, and on $H$
the evaluation is $-5$. Hence $\eta^2=25-36-1=-12$.
The projection formula for the component fundamental cycle gives
$\langle u^2x_i,[\widetilde X_p]\rangle=-12$, also when that
component has a self-identification: its normalization pushes its
fundamental class with degree one.

Finally $u$ is an actual pullback from downstairs. Therefore
\[
 \langle u^3,[\widetilde X_p]\rangle
 =\langle v^3,\pi_*[\widetilde X_p]\rangle
 =d\langle v^3,[X_p]\rangle=0\quad\text{in }\mathbb Z/d.
\]
This evaluates the otherwise undetermined finite-coefficient cube.
The exact support-function, wall-degree and intersection computations
are executed in \texttt{code/manifolds\_mu.py}.
\end{proof}
